% !TeX root = main.tex
% =====================================================================
%  preamble.tex --- packages, styling, title metadata
% =====================================================================

% --- encoding and language -------------------------------------------
\usepackage[T1]{fontenc}          % modern 8-bit font encoding (hyphenation, accents)
\usepackage[utf8]{inputenc}       % allow UTF-8 characters directly in the source
\usepackage[british]{babel}       % British English hyphenation and typographic conventions
\usepackage{csquotes}             % context-sensitive quotation marks (\enquote)

% --- maths -----------------------------------------------------------
\usepackage{amsmath}              % core AMS maths (align, aligned, equation, ...)
\usepackage{amsfonts}             % maths alphabets (\mathbb, \mathfrak, ...)
\usepackage{amssymb}              % extra maths symbols --- \Box (d'Alembertian) needs it
\usepackage{braket}               % Dirac bra-ket notation (\bra, \ket, \braket)
\usepackage{nicefrac}             % compact slanted in-line fractions (\nicefrac)
\usepackage{empheq}               % \empheq environment for boxing equations
\usepackage{environ}              % \NewEnviron --- captures an environment body in \BODY (used by keyeq)
\usepackage{mathrsfs}             % \mathscr --- used for the Hilbert space
\usepackage{siunitx}
%\usepackage{romannum}

% --- siunitx units setup ------------------------------------------------
\sisetup{per-mode=symbol-or-fraction} % 'slash' inline; normal fraction in a display
\sisetup{range-units = single} % setup of range units
\sisetup{range-phrase = {\text{~to~}}}   % system-wide rangephrase; \qtyrange[range-phrase=---]{0.7}{2.1}{\joule\per\cm^2} for per-instance changes

% --- floats, tables, graphics ----------------------------------------
\usepackage{graphicx}             % include images (\includegraphics)
\usepackage{float}                % [H] placement to pin floats exactly where written
\usepackage{multirow}             % table cells spanning several rows
\usepackage{tabularx}             % tables with auto-stretching X columns
\usepackage{wrapfig2}
\usepackage{caption}
\newcolumntype{L}{>{\raggedright\arraybackslash}X}

% --- layout ----------------------------------------------------------
\usepackage{placeins}             % \FloatBarrier to prevent floats from crossing a barrier 
\usepackage[nobottomtitles*]{titlesec} % customise section/paragraph heading styles
\usepackage{needspace}            % reserve N lines so a block is not split across a page
\usepackage[nottoc]{tocbibind}    % list bibliography/index in the ToC (not the ToC itself)
\usepackage{parskip}   	          % blank line between paragraphs, no first-line indent
\usepackage[left=1cm, top=1.5cm, bottom=2.5cm, right=5cm, marginparwidth=4.4cm, marginparsep=0.3cm, footskip=1cm, asymmetric]{geometry}
%\usepackage[margin=3.6cm, marginparwidth=3.1cm]{geometry}  % page margins; width of \todo notes
\usepackage{marginfix}

% --- colours  --------------------------------------------------------
\usepackage[dvipsnames]{xcolor}   % named colours (dvipsnames palette, e.g. Maroon)
\colorlet{linkcolor}{Maroon}      % single colour reused for all hyperlinks below

% --- hyperlinks ------------------------------------------------------
\usepackage{hyperref}             % clickable cross-references, URLs and ToC entries
\hypersetup{
	linkcolor  = linkcolor,
	citecolor  = linkcolor,
	urlcolor   = linkcolor,
	colorlinks = true,
	breaklinks = true
}
\usepackage[capitalize,nameinlink]{cleveref}  % smart references: \cref auto-inserts "Eq.", "Sec.", ...
\creflabelformat{equation}{#2#1#3}            % show equation refs as "3" rather than "(3)"

% --- title page metadata ---------------------------------------------
\title{%
	\vspace*{2cm}
	{\Huge Master Diploma Examination}\thanks{Translated into English with some modifications from an analogous text of \textit{Krokosz et al.}}\\[1cm]
	{\Large Revision Notes---Physics (Studies in English)}\\[0.4cm]
	{\large Faculty of Physics, University of Warsaw}\\
}

% --- comments --------------------------------------------------------
\usepackage[]{todonotes}                            % margin \todo notes used as revision flags; optional argument [disable] suppresses them for a clean, printable version (also change the margins then)

\newcommand{\todoper}[1]{\todo[backgroundcolor=red!25,linecolor=red!55]{#1}}       % red   --- persistent note, to be checked again later
\newcommand{\todoai}[1]{\todo[backgroundcolor=green!25,linecolor=green!55]{#1}}    % green --- rewritten and checked, comment from AI
\newcommand{\merr}{\textcolor{red}{\textbf{(!)}}} % marked mistakes (explanations in `%-comments`)
\newcommand{\todoans}[1]{\todo[backgroundcolor=blue!25,linecolor=blue!55]{#1}}  % blue  --- answer to a question I raised; for information only, not (yet) part of the text

% --- toc-link --------------------------------------------------------
\usepackage{fancyhdr}   % custom headers/footers (centred page number, ToC link at right)
\pagestyle{fancy}
\fancyhf{} % Clear all default header and footer fields
\renewcommand{\headrulewidth}{0pt} % Remove the top horizontal line
\fancyfoot[C]{\thepage}	% Place the page number in the centre
\fancyfoot[R]{\hyperref[toc]{TOC}} % Place the TOC link on the right-hand side ([L] for left). 

\author{???}
\date{\today}


% --- paragraph styling -----------------------------------------------
\titleformat{\paragraph}[hang]
    {\normalfont\normalsize} % Normal font size and standard weight
    {}                       % Empty label --- suppresses numbering
    {0pt}                    % Zero separation space since there is no number
    {\underline}             % Underlines the heading text
\titlespacing*{\section}{0pt}{3cm}{0.75cm}
\renewcommand{\bottomtitlespace}{.15\textheight}

% --- postulate environment (numbered) --------------------------------
% Roman numerals, counter runs document-wide, \label works inside so
% postulates can be cross-referenced with \cref{post:...}.
% \Needspace* (not \needspace) keeps the label, the italic body and a
% boxed equation together instead of splitting them over a page break.
\newcounter{postulate}
\renewcommand{\thepostulate}{\Roman{postulate}}
\newenvironment{postulate}[1]
    {\par\Needspace*{5\baselineskip}\medskip
     \refstepcounter{postulate}%
     \begingroup\leftskip=1.5em \rightskip=1.5em
     \noindent\textbf{Postulate \thepostulate\ ---\ #1.}\ \itshape\ignorespaces}
    {\par\endgroup\medskip}
\crefname{postulate}{postulate}{postulates}
\Crefname{postulate}{Postulate}{Postulates}



% --- contents lists (tocloft) ----------------------------------------
\usepackage{tocloft}

% ToC section entries: bold, with dot leaders (article supplies neither)
\renewcommand{\cftsecfont}{\bfseries}
\renewcommand{\cftsecpagefont}{\bfseries}
\renewcommand{\cftsecleader}{\bfseries\cftdotfill{\cftdotsep}}
\setlength{\cftsecindent}{0em}
\setlength{\cftsecnumwidth}{1.5em}

% List of Key Equations --- entries styled identically to ToC sections.
% No \numberline, so an entry is the display name and the page number alone.
\newlistof{keyequations}{loe}{List of Key Equations}
\renewcommand{\cftkeyequationsfont}{\bfseries}
\renewcommand{\cftkeyequationspagefont}{\bfseries}
\renewcommand{\cftkeyequationsleader}{\bfseries\cftdotfill{\cftdotsep}}
\setlength{\cftkeyequationsindent}{0em}
\setlength{\cftkeyequationsnumwidth}{0em}
\setlength{\cftbeforekeyequationsskip}{1.0em plus 1pt}

% Heading typeset as a \section*, exactly as article's \tableofcontents does,
% so \titlespacing and \titleformat apply to both lists identically.
\makeatletter
\renewcommand{\listofkeyequations}{%
  \section*{List of Key Equations}%
  \@starttoc{loe}}
\makeatother

\newcommand{\RomanNumeralCaps}[1]
    {\MakeUppercase{\romannumeral #1}}

% --- macro constants -------------------------------------------------
\newcommand{\Lagr}{\mathcal{L}}               % Lagrangian
\newcommand{\Ham}{\mathcal{H}}				        % Hamiltonian	
\newcommand{\Hilb}{\mathscr{H}}               % Hilbert space
\newcommand{\kB}{k_{B}}                       % Boltzmann constant
\newcommand{\muB}{\mu_{B}}                    % Bohr magneton
\newcommand{\const}{\mathrm{const.}}          % "constant"
\newcommand{\ev}[1]{\langle #1 \rangle}       % expectation / mean value
\newcommand{\comm}[2]{\left[#1,#2\right]}     % commutator
\newcommand{\acomm}[2]{\left\{#1,#2\right\}}  % anticommutator

% =====================================================================
%  Key equations --- capture once, replay in the appendix
% =====================================================================
% \begin{keyeq}{<label>}{<display name>} ... \end{keyeq}
%   * typesets the body boxed and numbered, exactly as the old
%     \begin{empheq}[box=\fbox]{equation} did,
%   * attaches \label{<label>},
%   * adds <display name> to the List of Key Equations,
%   * stores the body, the name, the label and the enclosing section so
%     \printkeyequations can reprint the lot in the appendix.
% There is therefore only ONE copy of each equation in the source and the
% appendix cannot drift from the text.
%
% Constraints on the body:
%   * no \label inside it (put the label in the first argument instead);
%     \printkeyequations gobbles stray ones, but the number would be wrong,
%   * anything legal inside `equation` is legal here --- aligned,
%     alignedat, gathered, arrays. A bare & or \\ at the top level is not.
%
% This block must sit AFTER hyperref, cleveref, titlesec and tocbibind,
% because it wraps whatever definition of \section they leave behind.

\makeatletter

% --- remember the title of the section currently being typeset --------
% \@currentlabelname is unusable here: a \caption between the heading and
% the equation would overwrite it with the caption text.
\newcommand{\keyeq@cursec}{}
\let\keyeq@oldsection\section
\renewcommand{\section}{\@ifstar{\keyeq@sectionstar}{\keyeq@section}}
\newcommand{\keyeq@sectionstar}[1]{\keyeq@oldsection*{#1}}
\newcommand{\keyeq@section}[2][\relax]{%
  \gdef\keyeq@cursec{#2}%
  \ifx\relax#1\relax\keyeq@oldsection{#2}\else\keyeq@oldsection[#1]{#2}\fi}

\newcounter{keyeqno}

\NewEnviron{keyeq}[2]{%
  \stepcounter{keyeqno}%
  \expandafter\global\expandafter\let\csname keyeq@body@\thekeyeqno\endcsname\BODY
  \expandafter\gdef\csname keyeq@name@\thekeyeqno\endcsname{#2}%
  \expandafter\xdef\csname keyeq@lab@\thekeyeqno\endcsname{#1}%
  \expandafter\xdef\csname keyeq@secno@\thekeyeqno\endcsname{\thesection}%
  \expandafter\global\expandafter\let\csname keyeq@sectitle@\thekeyeqno\endcsname\keyeq@cursec
  \begin{empheq}[box=\fbox]{equation}\label{#1}
    \BODY
  \end{empheq}%
  \phantomsection
  \addcontentsline{loe}{keyequations}{#2}%
}

% --- reprint every captured equation, grouped by section --------------
\newcommand{\printkeyequations}{%
  \begingroup
  \let\label\@gobble                     % never re-issue a label on replay
  \def\keyeq@prevsecno{\relax}%
  \count@=0\relax
  \loop\ifnum\count@<\value{keyeqno}\relax
    \advance\count@ by 1\relax
    \edef\keyeq@i{\the\count@}%
    \edef\keyeq@thissecno{\csname keyeq@secno@\keyeq@i\endcsname}%
    \ifx\keyeq@thissecno\keyeq@prevsecno\else
      \Needspace*{6\baselineskip}%
      \bigskip\noindent{\Large\bfseries\csname keyeq@sectitle@\keyeq@i\endcsname}\par
      \vspace{6pt}%
      \let\keyeq@prevsecno\keyeq@thissecno
    \fi
    \Needspace*{6\baselineskip}%
    \par\noindent\hspace*{3em}\textbf{\csname keyeq@name@\keyeq@i\endcsname}\hfill
      \cref{\csname keyeq@lab@\keyeq@i\endcsname},
      p.~\pageref{\csname keyeq@lab@\keyeq@i\endcsname}\par
    \vspace{-0.5\baselineskip}%
    \begin{equation*}
      \csname keyeq@body@\keyeq@i\endcsname
    \end{equation*}
  \repeat
  \endgroup
}

\makeatother


% =====================================================================
%  House style --- full version in STYLE.md
% =====================================================================
% \emph{} for defined terms and for emphasis in prose.
% \textbf{} ONLY as a label at the head of a list item. Never mid-sentence.
% \textit{} ONLY for a law or theorem quoted as a self-contained statement.
% \nicefrac{}{} for every inline fraction, differential operators included.
%   \frac and \tfrac in display maths only.
% bmatrix (square brackets) for every vector and matrix, never pmatrix.
% Inline matrices: \left[\begin{smallmatrix}...\end{smallmatrix}\right].
% Vector-valued operators carry hat and arrow: \hat{\vec{L}}. Components: \hat{L}_{z}.
% Differentials are a plain italic d, not \mathrm{d}. Inexact ones are \delta Q.
% f is frequency, never \nu. \varepsilon_{0}, never \epsilon_{0}.
% Equations that must be memorised: \begin{keyeq}{eq:label}{Display name}.
%   Boxed, numbered, listed and replayed in the appendix. No \label inside the body.
% No semicolons in prose. British spelling. \enquote{} for quotation marks.
% Em-dash as --- with no surrounding spaces. Three exceptions, all spaced:
%   a symbol gloss after a display, a label gloss at the head of a list item,
%   and section titles, which reproduce the exam question as written.
% One paragraph per source line, blank line between paragraphs.